(1) 各向同性均匀介电体的慢极化 \(\boldsymbol P\) 满足 \(\tau\dot{\boldsymbol P}+\boldsymbol P=\varepsilon_0\Delta\varepsilon\,\boldsymbol E(t)\)，其中 \(\tau>0,\Delta\varepsilon>0\)，另有瞬时高频响应 \(\varepsilon_\infty>0\)。从 \(t=0\) 起施加阶跃均匀电场 \(E_0\hat x\)，初始慢极化为零。求时间响应与因果核。
(2) 对 \(E(t)=\operatorname{Re}[E_0e^{-i\omega t}]\) 求复介电常数 \(\varepsilon^*(\omega)=\varepsilon'(\omega)+i\varepsilon''(\omega)\)，给出实部、虚部及静态与高频极限。
(3) 计算单位体积每周期电场对极化所作的净功，求介电损耗峰的位置和峰高，解释被动介质为何必须满足 \(\varepsilon''\ge0\)（本题时间约定）。
(4) 独立验证 \(\Delta\varepsilon=(2/\pi)\int_0^\infty\varepsilon''(\omega)\,d\omega/\omega\) 的因果色散和则，并讨论若 \(\tau\to0\) 或 \(\Delta\varepsilon\to0\) 时的响应。
